\documentclass[../report_main.tex]{subfiles}
\begin{document}

\section{Introduction}
A robust database application requires not only standard \gls{CRUD} queries but also advanced database objects such as Views, Functions, Triggers, and Indexes to ensure data integrity, optimize performance, and handle complex business logic natively within the database engine. This chapter details 47 distinct SQL statements categorized into Complex Application Queries (extracted from the backend API), User-Defined Functions, Triggers, and Indexing Strategies.

\section{Authentication \& Profile Queries}

\textbf{Query 1: Check login information}
\begin{lstlisting}[language=SQL]
SELECT * FROM accounts 
WHERE email = 'teacher1@email.com' AND is_active = true;
-- Password verification is handled at the application level using a Python function
\end{lstlisting}

\textbf{Query 2: Retrieve full user profile with role-specific details}
\begin{lstlisting}[language=SQL]
SELECT a.id, a.email, a.role, a.name, a.avatar_url,
       t.organization, s.school
FROM accounts a
LEFT JOIN teachers t ON a.id = t.id AND a.role = 'teacher'
LEFT JOIN students s ON a.id = s.id AND a.role = 'student'
WHERE a.id = 12345;
\end{lstlisting}

\textbf{Query 3: Create a new user account}
\begin{lstlisting}[language=SQL]
INSERT INTO accounts (email, password, role, name) 
VALUES ('teacher@email.com', 'password', 'teacher', 'Teacher s name')
RETURNING id;
\end{lstlisting}

\textbf{Query 4: Create corresponding user profile}
\begin{lstlisting}[language=SQL]
INSERT INTO teachers (id, organization) 
VALUES (123456, 'SoICT-HUST');

INSERT INTO students (id, school)
VALUES (123456, 'Hanoi University of Science and Technology');
\end{lstlisting}

\textbf{Query 5: Update user profile}
\begin{lstlisting}[language=SQL]
UPDATE accounts
SET name = 'Teacher s name' WHERE id = 123456
RETURNING new.*;
\end{lstlisting}

\textbf{Query 6: Change account password}
\begin{lstlisting}[language=SQL]
UPDATE accounts
SET password = 'new_hashed_password_string' 
WHERE id = 123456;
-- Password verification is handled at the application level using a Python function
\end{lstlisting}

\section{Data Upload \& Modification Queries}

\textbf{Query 7: Upload a new question}
\begin{lstlisting}[language=SQL]
INSERT INTO questions (
    teacher_id, public_id, subject, grade, parent_id,
    question_type, layout_type, content, solution,
    chapter, lesson, complexity, is_shufflable
)
VALUES (
    123456, '1b95845a...', 'Physics', '12', NULL,
    'mc', 'normal', 'Question Content...', 
    'Detailed Solution...', 'Chapter 1', 'Lesson 1', 1, TRUE
)
RETURNING id;
\end{lstlisting}

\textbf{Query 8: Upload question's image}
\begin{lstlisting}[language=SQL]
INSERT INTO q_images (question_id, storage_path, img_type, img_scale) 
VALUES (123456, '/static/img.png', 'png', 1.0);
\end{lstlisting}

\textbf{Query 9: Upload question's options}
\begin{lstlisting}[language=SQL]
INSERT INTO q_choice_details (question_id, content, is_correct, order_index, is_shufflable) 
VALUES 
    (123456, 'Option A', TRUE, 1, TRUE),
    (123456, 'Option B', FALSE, 2, TRUE),
    (123456, 'Option C', FALSE, 3, TRUE),
    (123456, 'Option D', FALSE, 4, FALSE);
\end{lstlisting}

\textbf{Query 10: Update question content}
\begin{lstlisting}[language=SQL]
UPDATE questions 
SET content = 'New content...', complexity = 4 
WHERE id = 123456;
\end{lstlisting}

\textbf{Query 11: Update question options}
\begin{lstlisting}[language=SQL]
UPDATE q_choice_details 
SET content = 'New Option' 
WHERE id = 11223344 
AND question_id = 123456;
\end{lstlisting}

\textbf{Query 12: Update question's image scaling}
\begin{lstlisting}[language=SQL]
UPDATE q_images 
SET img_scale = 0.8 
WHERE question_id = 123456;
\end{lstlisting}

\textbf{Query 13: Soft delete question}
\begin{lstlisting}[language=SQL]
UPDATE questions 
SET deleted_at = NOW() 
WHERE id = 111 
AND teacher_id = 81; 
-- Prevent teacher A try to delete teacher B question
\end{lstlisting}

\textbf{Query 14: Restore soft-deleted question}
\begin{lstlisting}[language=SQL]
UPDATE questions 
SET deleted_at = NULL 
WHERE id = 111 
AND teacher_id = 81;
\end{lstlisting}

\textbf{Query 15: Permanently delete question}
\begin{lstlisting}[language=SQL]
DELETE FROM questions 
WHERE id = 11539 
AND teacher_id = 81;
\end{lstlisting}


\section{Question Bank Managing Queries}

\textbf{Query 16: Count total number of questions based on filtering criteria}
\begin{lstlisting}[language=SQL]
SELECT COUNT(*) as total 
FROM questions q 
WHERE q.deleted_at IS NULL 
AND q.parent_id IS NULL 
AND q.teacher_id = 81 
AND q.subject = 'Physics' 
AND q.grade = 12 
AND (q.content ILIKE '%BJT transistor%' OR EXISTS 
    (SELECT 1 
    FROM questions child 
    WHERE child.parent_id = q.id 
    AND child.content ILIKE '%BJT transistor%')
);
\end{lstlisting}

\textbf{Query 17: Query first page's contents (Pagination)}
\begin{lstlisting}[language=SQL]
SELECT q.*, t.name AS teacher_name
FROM questions q
LEFT JOIN accounts t ON q.teacher_id = t.id
WHERE q.deleted_at IS NULL 
AND q.parent_id IS NULL 
AND q.teacher_id = 81
ORDER BY q.id DESC 
LIMIT 20 OFFSET 0;
\end{lstlisting}

\textbf{Note:} Queries 18, 19, and 20 demonstrate a batch query optimization technique (using the \texttt{IN (...)} clause) to resolve the $N+1$ query problem. Instead of executing separate database calls to fetch images and options for each individual question, the system aggregates all question IDs from the paginated result and fetches their respective relations in a single database round-trip, significantly reducing latency and server load.

\textbf{Query 18: Query questions image}
\begin{lstlisting}[language=SQL]
SELECT question_id, storage_path, img_scale 
FROM q_images 
WHERE question_id IN (1000, 1010, 1005);
\end{lstlisting}

\textbf{Query 19: Query questions options}
\begin{lstlisting}[language=SQL]
SELECT id, question_id, content, is_correct 
FROM q_choice_details 
WHERE question_id IN (4, 5);
\end{lstlisting}

\textbf{Query 20: Query following questions (for group/stimulus questions)}
\begin{lstlisting}[language=SQL]
SELECT id, parent_id, question_type, content, solution, complexity 
FROM questions 
WHERE parent_id IN (1431, 481, 502) 
ORDER BY id ASC;
\end{lstlisting}


\section{Contest Managing Queries}

\textbf{Query 21: Retrieve teacher's created contests}
\begin{lstlisting}[language=SQL]
SELECT ct.*, cl.class_name,
       COALESCE(SUM(CASE WHEN q.question_type != 'st' THEN 1 ELSE 0 END), 0) as question_count
FROM contests ct
LEFT JOIN classes cl ON cl.id = ct.class_id
LEFT JOIN contests_questions cq ON cq.contest_id = ct.id
LEFT JOIN questions q ON q.id = cq.question_id
WHERE ct.teacher_id = 4
GROUP BY ct.id, cl.class_name
ORDER BY ct.id DESC;
\end{lstlisting}

\textbf{Query 22: Retrieve student's enrolled contests}
\begin{lstlisting}[language=SQL]
SELECT ct.*, cl.class_name
FROM contests ct
JOIN students_classes sc ON sc.class_id = ct.class_id
JOIN classes cl ON cl.id = ct.class_id
WHERE sc.student_id = 10000 
AND ct.status = 'active';
\end{lstlisting}

\textbf{Query 23: Create a new contest}
\begin{lstlisting}[language=SQL]
INSERT INTO contests (class_id, teacher_id, title, time_limit, scoring_config, status) 
VALUES (5, 4, 'Midterm Exam', 45, '{"mc_weight": 1.0}', 'active') 
RETURNING id, public_id;
\end{lstlisting}

\textbf{Query 24: Insert questions into contest}
\begin{lstlisting}[language=SQL]
INSERT INTO contests_questions (contest_id, question_id, original_order, point_weight) 
VALUES (20, 150, 1, 2.0);
\end{lstlisting}

\textbf{Query 25: Open/Close a contest}
\begin{lstlisting}[language=SQL]
UPDATE contests 
SET status = 'inactive' -- Or active/deleted (Soft delete)
WHERE id = 103 
AND teacher_id = 4;
\end{lstlisting}

\textbf{Query 26: Permanently delete a contest}
\begin{lstlisting}[language=SQL]
DELETE FROM contests 
WHERE id = 103 
AND teacher_id = 4;
\end{lstlisting}

\textbf{Query 27: Create new contest result (Student starts attempt)}
\begin{lstlisting}[language=SQL]
INSERT INTO contest_results (student_id, contest_id, guest_name, start_time) 
VALUES (10000, 20, NULL, NOW()) 
RETURNING id;
\end{lstlisting}

\textbf{Query 28: Create new student submission (Student submits a question)}
\begin{lstlisting}[language=SQL]
INSERT INTO student_option_submissions (
    contest_result_id, question_id, student_choice, option_display_order, is_correct, earned_point
) 
VALUES (55, 205, 'assignment operator =', '', TRUE, 1.0);
\end{lstlisting}

\textbf{Query 29: Update contest result (Student submits the attempt)}
\begin{lstlisting}[language=SQL]
UPDATE contest_results
SET end_time = NOW(), total_score = 8.5, count_wrong_answers = 2
WHERE id = 55;
\end{lstlisting}

\textbf{Query 30: Show a contest's results (Teacher view)}
\begin{lstlisting}[language=SQL]
SELECT cr.id AS result_id, cr.start_time, cr.end_time, cr.total_score, cr.count_wrong_answers,
       s.name AS student_name, cr.guest_name
FROM contest_results cr
LEFT JOIN accounts s ON s.id = cr.student_id
WHERE cr.contest_id = 20
ORDER BY cr.id DESC;
\end{lstlisting}

\textbf{Query 31: Show student's submission details (Student review)}
\begin{lstlisting}[language=SQL]
SELECT sos.*, q.content, q.solution
FROM student_option_submissions sos
JOIN questions q ON q.id = sos.question_id
WHERE sos.contest_result_id = 55;
\end{lstlisting}


\section{Class Managing Queries}

\textbf{Query 32: Retrieve class lists for Teacher}
\begin{lstlisting}[language=SQL]
SELECT c.*, COUNT(sc.student_id) AS student_count, COUNT(ct.id) AS contest_count
FROM classes c
LEFT JOIN students_classes sc ON sc.class_id = c.id
LEFT JOIN contests ct ON ct.class_id = c.id
WHERE c.teacher_id = 4 
GROUP BY c.id 
ORDER BY c.create_at DESC;
\end{lstlisting}

\textbf{Query 33: Retrieve class lists for Student}
\begin{lstlisting}[language=SQL]
SELECT c.*, t.name as teacher_name, COUNT(sc2.student_id) AS student_count
FROM classes c
JOIN students_classes sc ON sc.class_id = c.id AND sc.student_id = 10000
LEFT JOIN accounts t ON t.id = c.teacher_id
LEFT JOIN students_classes sc2 ON sc2.class_id = c.id
GROUP BY c.id, t.name 
ORDER BY c.create_at DESC;
\end{lstlisting}

\textbf{Query 34: Create a new class}
\begin{lstlisting}[language=SQL]
INSERT INTO classes (teacher_id, class_name, description) 
VALUES (100, 'Advanced JavaFX Programming', 'Practical class...') 
RETURNING id, public_id;
\end{lstlisting}

\textbf{Query 35: Delete a class}
\begin{lstlisting}[language=SQL]
DELETE FROM classes 
WHERE id = 5 
AND teacher_id = 100;
\end{lstlisting}

\textbf{Query 36: Query class's students}
\begin{lstlisting}[language=SQL]
SELECT a.id, a.name, a.email, sc.create_at as joined_at
FROM students_classes sc
JOIN accounts a ON a.id = sc.student_id
WHERE sc.class_id = 5;
\end{lstlisting}

\textbf{Query 37: Query class's contests}
\begin{lstlisting}[language=SQL]
SELECT id, title, status, time_limit 
FROM contests 
WHERE class_id = 5 
ORDER BY id DESC;
\end{lstlisting}

\textbf{Query 38: Search for students to add}
\begin{lstlisting}[language=SQL]
SELECT a.id, a.name, a.email
FROM accounts a
WHERE a.role = 'student'
AND (CAST(a.id AS TEXT) LIKE '%son%' 
OR a.email ILIKE '%son%' 
OR a.name ILIKE '%son%')
LIMIT 10;
\end{lstlisting}

\textbf{Query 39: Add student to class}
\begin{lstlisting}[language=SQL]
INSERT INTO students_classes (student_id, class_id) 
VALUES (10000, 5);
\end{lstlisting}

\textbf{Query 40: Remove student from class}
\begin{lstlisting}[language=SQL]
DELETE FROM students_classes 
WHERE student_id = 10000 
AND class_id = 5;
\end{lstlisting}


\section{Database Functions and Triggers}

\textbf{Query 41: Auto-sync teacher's question count}
\begin{lstlisting}[language=SQL]
ALTER TABLE teachers ADD COLUMN question_count INT DEFAULT 0;

UPDATE teachers t
SET question_count = COALESCE((
    SELECT COUNT(*) 
    FROM questions q 
    WHERE q.teacher_id = t.id AND q.parent_id IS NULL AND q.deleted_at IS NULL
), 0);

CREATE OR REPLACE FUNCTION trg_sync_teacher_question_count()
RETURNS TRIGGER AS $$
BEGIN
    IF (TG_OP = 'INSERT') AND NEW.parent_id IS NULL 
    THEN
        UPDATE teachers 
        SET question_count = question_count + 1 
        WHERE id = NEW.teacher_id;
        RETURN NEW;
    ELSIF (TG_OP = 'DELETE') AND OLD.parent_id IS NULL 
    THEN
        UPDATE teachers 
        SET question_count = question_count - 1 
        WHERE id = OLD.teacher_id;
        RETURN OLD;
    ELSIF (TG_OP = 'UPDATE') AND NEW.parent_id IS NULL 
    THEN
        IF (NEW.deleted_at IS NOT NULL AND OLD.deleted_at IS NULL) 
        THEN
            UPDATE teachers 
            SET question_count = question_count - 1 
            WHERE id = NEW.teacher_id;
        ELSIF (NEW.deleted_at IS NULL AND OLD.deleted_at IS NOT NULL) 
        THEN
            UPDATE teachers 
            SET question_count = question_count + 1 
            WHERE id = NEW.teacher_id;
        END IF;
        RETURN NEW;
    END IF;
    RETURN NULL;
END;
$$ LANGUAGE plpgsql;

CREATE TRIGGER sync_question_count
AFTER INSERT OR UPDATE OR DELETE ON questions
FOR EACH ROW EXECUTE FUNCTION trg_sync_teacher_question_count();
\end{lstlisting}

\textbf{Query 42: Auto-sync class's student count}
\begin{lstlisting}[language=SQL]
ALTER TABLE classes ADD COLUMN student_count INT DEFAULT 0;

UPDATE classes c
SET student_count = COALESCE((
    SELECT COUNT(*) 
    FROM students_classes sc 
    WHERE sc.class_id = c.id
), 0);

CREATE OR REPLACE FUNCTION trg_sync_class_student_count()
RETURNS TRIGGER AS $$
BEGIN
    IF (TG_OP = 'INSERT') 
    THEN
        UPDATE classes 
        SET student_count = student_count + 1 
        WHERE id = NEW.class_id;
        RETURN NEW;
    ELSIF (TG_OP = 'DELETE') 
    THEN
        UPDATE classes 
        SET student_count = student_count - 1 
        WHERE id = OLD.class_id;
        RETURN OLD;
    END IF;
    RETURN NULL;
END;
$$ LANGUAGE plpgsql;

CREATE TRIGGER sync_student_count
AFTER INSERT OR DELETE ON students_classes
FOR EACH ROW EXECUTE FUNCTION trg_sync_class_student_count();
\end{lstlisting}

\textbf{Query 43: Auto-sync class's contest count}
\begin{lstlisting}[language=SQL]
ALTER TABLE classes ADD COLUMN contest_count INT DEFAULT 0;

UPDATE classes c
SET contest_count = COALESCE((
    SELECT COUNT(*) 
    FROM contests ct 
    WHERE ct.class_id = c.id
), 0);

CREATE OR REPLACE FUNCTION trg_sync_class_contest_count()
RETURNS TRIGGER AS $$
BEGIN
    IF (TG_OP = 'INSERT') 
    THEN
        UPDATE classes 
        SET contest_count = contest_count + 1 
        WHERE id = NEW.class_id;
        RETURN NEW;
    ELSIF (TG_OP = 'DELETE') 
    THEN
        UPDATE classes 
        SET contest_count = contest_count - 1 
        WHERE id = OLD.class_id;
        RETURN OLD;
    END IF;
    RETURN NULL;
END;
$$ LANGUAGE plpgsql;

CREATE TRIGGER sync_contest_count
AFTER INSERT OR DELETE ON contests
FOR EACH ROW EXECUTE FUNCTION trg_sync_class_contest_count();
\end{lstlisting}


\section{Database Indexes (Optimization Strategy)}

\textbf{Query 44: Authentication Optimization}
\begin{lstlisting}[language=SQL]
CREATE UNIQUE INDEX idx_accounts_email ON accounts(email);
\end{lstlisting}

\textbf{Query 45: Question Bank Query Optimization}
\begin{lstlisting}[language=SQL]
CREATE INDEX idx_questions_search ON questions(teacher_id, subject, grade) 
WHERE deleted_at IS NULL;

CREATE INDEX idx_questions_parent_id ON questions(parent_id) 
WHERE parent_id IS NOT NULL;

CREATE INDEX idx_q_choice_details_question_id ON q_choice_details(question_id);

CREATE INDEX idx_q_images_question_id ON q_images(question_id);
\end{lstlisting}

\textbf{Query 46: Contest and Submission Optimization}
\begin{lstlisting}[language=SQL]
CREATE INDEX idx_contests_teacher_id ON contests(teacher_id);

CREATE INDEX idx_contests_questions_contest_id ON contests_questions(contest_id);

CREATE INDEX idx_contest_results_contest_id ON contest_results(contest_id);
\end{lstlisting}

\textbf{Query 47: Class Management Optimization}
\begin{lstlisting}[language=SQL]
CREATE INDEX idx_sc_student_id ON students_classes(student_id);

CREATE INDEX idx_sc_class_id ON students_classes(class_id);

CREATE UNIQUE INDEX idx_classes_public_id ON classes(public_id);
\end{lstlisting}
\textbf{Query 48: Dashboard Query Optimization}
\begin{lstlisting}[language=SQL]
CREATE INDEX idx_questions_dashboard ON questions(teacher_id) 
WHERE deleted_at IS NULL AND parent_id IS NULL;

CREATE INDEX idx_classes_teacher_id ON classes(teacher_id);

CREATE INDEX idx_contests_class_id ON contests(class_id);
\end{lstlisting}

\end{document}